A redução necessária é dada pela equação:

\begin{equation}
  I_{\text{nec}} = \frac{n_{\text{motor}}}{n_{\text{tambor}}}
  \label{eq:reducao_necessaria}
\end{equation}

Onde:

\begin{itemize}[leftmargin=1.25cm]
  \item $I_{\text{nec}}$: redução necessária.
  \item $n_{\text{motor}}$: rotação do motor em $\si{\rpm}$.
  \item $n_{\text{tambor}}$: rotação do tambor em $\si{\rpm}$.
\end{itemize}

\vspace{1em}

A velocidade do cabo é dada pela equação:

\begin{equation}
  V_{\text{cabo}} = \frac{n_c}{2} \cdot V_L
  \label{eq:velocidade_do_cabo}
\end{equation}

Onde:

\begin{itemize}[leftmargin=1.25cm]
  \item $V_{\text{cabo}}$: velocidade do cabo em $\si{\meter/\second}$.
  \item $n_c$: número de cabos.
  \item $V_L$: velocidade de elevação em $\si{\meter/\second}$.
\end{itemize}

\vspace{1em}

Substituindo os valores na equação \eqref{eq:velocidade_do_cabo}:

\begin{align*}
  V_{\text{cabo}} = \frac{4}{2} \cdot 0,1 \\
  V_{\text{cabo}} = \boxed{\SI{0,2}{\meter/\second}}
\end{align*}

A velocidade angular do tambor é dada pela equação:

\begin{equation}
  \omega_{\text{tambor}} = \frac{V_{\text{cabo}}}{R_t}
  \label{eq:velocidade_angular_do_tambor}
\end{equation}

Onde:

\begin{itemize}[leftmargin=1.25cm]
  \item $\omega_{\text{tambor}}$: velocidade angular do tambor em $\si{\radian/\second}$.
  \item $V_{\text{cabo}}$: velocidade do cabo em $\si{\meter/\second}$.
  \item $R_t$: raio do tambor em $\si{\meter}$.
\end{itemize}

\vspace{1em}

Substituindo os valores na equação \eqref{eq:velocidade_angular_do_tambor}:

\begin{align*}
  \omega_{\text{tambor}} = \frac{0,2}{\frac{355,6}{2} \cdot 10^{-3}} \\
  \omega_{\text{tambor}} = \boxed{\SI{1,13}{\radian/\second}}
\end{align*}

\begin{align*}
  n_{\text{tambor}} = \frac{60}{2 \pi} \cdot 1,13 \\
  n_{\text{tambor}} = \boxed{\SI{10,79}{\rpm}}
\end{align*}

Substituindo os valores na equação \eqref{eq:reducao_necessaria}:

\begin{align*}
  I_{\text{nec}} = \frac{1170}{10,79} \\
  I_{\text{nec}} = \boxed{108,43}
\end{align*}
